\begin{claim}{D4.5}{Proved}{The full twisted cohomology and its jump loci}
For $z\notin\Z$, let $\mathcal R_z$ be the set of regular base
points satisfying (D.15), counted as points, not as marked integer
tuples. It is finite. Put $N_z=|\mathcal R_z|$. Define
$\ell_1(z)=0$ off $\Lambda_1$, and
$\ell_1(z)=1+\epsilon_n$ if
$z=n+k(3\beta_1+4-2\rho)$, $k\ne0$.
Define $\ell_2(z)=1$ if
$z=n+k(2\beta_2+3-3i)$, $k\ne0$, with
$n\equiv1,2\pmod4$, and $\ell_2(z)=0$ otherwise.
Then, setting $J_z=2N_z+\ell_1(z)+\ell_2(z)$,
\[
 \boxed{(h^0,h^1,h^2,h^3)(T_X\otimes M_z)=(0,J_z,1+J_z,0).} \tag{D.25}
\]
More precisely, with $\mathfrak m_t$ the maximal ideal at $t$,
\[
 R^1f_*(T_X\otimes M_z)=
   \bigoplus_{t\in\mathcal R_z}\C_t^2
      \oplus\C_{p_1}^{\ell_1(z)}\oplus\C_{p_2}^{\ell_2(z)},
\]
\[
 R^2f_*(T_X\otimes M_z)=\C_{p_0}\oplus
   \bigoplus_{t\in\mathcal R_z}\OO_{B,t}/\mathfrak m_t^2
      \oplus\C_{p_1}^{\ell_1(z)}\oplus\C_{p_2}^{\ell_2(z)}. \tag{D.26}
\]
Together with D4.1 these formulas cover every element of
$\operatorname{Pic}(X)=\C$. The actual cohomology jump loci are
locally discrete subsets of this complex line.
\end{claim}
\begin{proof}
The support assertions and all stalks were proved in D4.2--D4.4.
The direct image in degree zero is zero for a nonintegral twist:
it has generic rank zero and is torsion-free, since multiplication
by a nonzero base function is injective on the bundle on the smooth
integral total space. The support of either higher coherent direct
image is a proper closed analytic subset of the compact curve,
and hence is finite. This proves finiteness of $\mathcal R_z$ as
well as the absence of an accumulation at the special values.
If two integer tuples at a marked fibre represented the same
trivialized line, their difference would be an integral real class
in $H^{1,0}$, which is zero. Thus there is no multiplicity from
duplicate tuples. Formula (D.26) follows by the decomposition of a
finite-support sheaf into its local modules. Its global sections
have dimension equal to its length and its $H^1$ is zero, so Leray
gives (D.25).

To clarify the word generic, the global closed form $G_\gamma$ in
C3.8 defines a holomorphic line bundle on $X\times\C$ by the
integrable operators $\bar\partial-2\pi i zG_\gamma$. Proper
coherent base change for the projection to $\C$ makes each
cohomology-dimension threshold an analytic subset. A real noninteger
$z$ has no resonance by (D.17) and (D.16), so the generic values are
attained. All possible exceptional values were shown to form a
countable set. Thus each actual threshold exceeding a generic
dimension is a proper analytic subset of a complex curve and is
locally discrete. Some indexed resonances, namely the two invisible
order-four congruences, produce no jump; (D.25) distinguishes them.

This is an explicit countable description of the jump loci in the
Picard coordinate, rather than a generic assertion. It can also be
indexed algebraically on the elliptic surface. If a regular fibre
has torsion section $P$ of exact order $r_t$, choose its marked
relation $q_t\tau+r_t\mu=p_t$. Its full resonance lattice is
\[
 \Z+\bigl(6q_t\mu(t)+r_t\beta_c(t)\bigr)\Z.
\]
Every regular resonance arises this way. The underlying values
$t$ are exactly the non-special roots of the division polynomials
$\psi_N(P)$, $N\ge2$. The nonintegral lattice points contribute two
to $J_z$ for each distinct such fibre. The finite-point lattices
have the congruence weights already stated. The script
\texttt{d\_division\_polynomials.py} evaluates the polynomials
exactly for $2\le N\le12$; for example the first nonconstant
regular factors are $5t-32$, $89t^2+4736t-4096$ and
$11t^2-6016t+8192$. These computations specify an executable
enumeration of the resonance loci without inventing decimal values
for the unspecified parameter $c$.

As an independent index check, $c_1(X)=c_2(X)=c_1(M)=0$ in integral
cohomology, so the degree-six Riemann--Roch integrand for the actual
twisted tangent bundle is $c_3/2$. Its integral is one. Both (D.13)
and (D.25) have alternating sum one. The index is used here to
verify the completed sheaf calculation, not to select its lengths.
\end{proof}
